% 01_introduction.tex
%
% Purpose:
%   Introduce the line-defect waveguide problem, the constructive trace-subspace
%   method, and the relation to existing quasiperiodic BIE and periodic
%   waveguide boundary-condition literature.
%
% Main algorithm:
%   This section is expository.  It states the computational objective and
%   positions the formulation before the mathematical details are introduced.
%
% Based on:
%   sections/intro.tex from the first report draft and pre/pre1/pre1.tex.
%
% Main changes:
%   Keeps the motivation and literature discussion at the top level while
%   moving detailed geometry and trace notation to the preliminaries.
%
% Numerical goal:
%   Clarify what the report is meant to support: scattering solves, homogeneous
%   singular-value scans, and future numerical validation.

\section{Introduction}

This report develops a constructive numerical realization for periodic
line-defect waveguide scattering and eigenvalue-type calculations.  The model is
a scalar Helmholtz transmission problem in a two-dimensional geometry that is
quasiperiodic in the transverse direction and open through periodic half-leads
in the longitudinal direction.  The two main workflows are forced lead
scattering, where incoming Bloch data are prescribed from one or both leads, and
homogeneous eigenvalue or transmission-eigenvalue-type searches, where a
nontrivial admissible null vector is sought.

The formulation combines several standard numerical ingredients in a compact
port framework: boundary integral equations on material interfaces, artificial
quasiperiodic ports, Rayleigh coordinates for wall Cauchy traces, one-cell
scattering computations in the periodic leads, Bloch-mode selection of
incoming/outgoing trace subspaces, and a center-cell matching system.  Its main
point is not to construct a new abstract Dirichlet-to-Neumann theory for
periodic half-guides.  Instead, it builds the finite-dimensional Cauchy trace
subspaces needed by the current line-defect geometry and imposes those
subspaces directly in the global matrix.

The distinction between coordinates and physics is central.  Rayleigh modes form
a convenient Fourier basis for traces on a vertical port, but an arbitrary
Rayleigh Cauchy vector need not be the trace of an admissible half-lead solution.
Physical incoming and outgoing data are identified by a lead-cell Bloch or
scattering computation, followed by a mode-selection rule.  Once the admissible
subspace has been selected, any numerically well-conditioned basis of that same
subspace may be used in the coupled system.

\subsection{Relation to existing methods}

The present formulation is closely related to several established strands of
work.  Boundary integral methods for quasiperiodic band-structure calculations,
including the numerical difficulties associated with quasiperiodic Green
functions and empty-cell resonances, are represented by the work of Barnett and
Greengard \cite{BarnettGreengard2010}.  Exact boundary conditions for locally
perturbed periodic waveguides are developed by Joly, Li, and Fliss
\cite{JolyLiFliss2006}, where half-guide exterior conditions are described
through periodic-guide Dirichlet-to-Neumann and related operators.  Coatl\'even
\cite{Coatleven2012} treats Helmholtz problems in periodic media with line
defects using a combination of Floquet--Bloch transforms,
Lippmann--Schwinger ideas, and Dirichlet-to-Neumann maps.

The present formulation differs from these approaches in emphasis.  It is a
compact-port finite-dimensional trace-space realization tailored to the current
boundary integral implementation.  It should not be read as replacing the full
functional-analytic half-guide theory.  Rather, it records how Rayleigh wall
coordinates, one-cell Bloch/scattering calculations, and center-cell boundary
integral equations are assembled into a practical matrix formulation whose
truncation and mode-selection choices can be tested numerically.
